\section{Interior drift and Gaussian fluctuations}
\label{sec:limits}

Throughout this section, $0<p<1$, $q=1-p$, and $a=pq$.  At zero tilt,
$P_p:=Q_p(1)$ is the transition matrix of the lumped projective chain.

\begin{lemma}[Stationary law of the lumped projective chain]
\label{lem:stationary}
The unique stationary row vector of $P_p$ is
\begin{equation}
 \label{eq:stationary}
 \boxed{\quad
 \pi_N=\frac{p}{1+p},\qquad
 \pi_Z=\frac{1-a}{2+a},\qquad
 \pi_P=\frac{q}{1+q}.
 \quad}
\end{equation}
Its stationary mean reward is
\begin{equation}
 \label{eq:drift}
 \boxed{\quad \mu_p=\frac{3a}{2+a}>0.\quad}
\end{equation}
\end{lemma}

\begin{proof}
Every transition into $N$ uses $A$, except the transition from $N$ itself.
Stationarity therefore gives $\pi_N=p(1-\pi_N)$ and hence
$\pi_N=p/(1+p)$.  Likewise, $\pi_P=q(1-\pi_P)$, so
$\pi_P=q/(1+q)$.  The remaining mass is
\[
 1-\frac{p}{1+p}-\frac{q}{1+q}=\frac{1-a}{2+a}.
\]
Only $N\xrightarrow{A}Z$ and $P\xrightarrow{B}Z$ carry reward $-1$.
All other transitions carry reward $+1$.  Thus
\begin{align*}
 \mu_p
 &=1-2(p\pi_N+q\pi_P)\\
 &=1-2\left(\frac{p^2}{1+p}+\frac{q^2}{1+q}\right)
 =1-\frac{2(1-a)}{2+a}=\frac{3a}{2+a}.
\end{align*}
\end{proof}

\begin{theorem}[SLLN and explicit CLT]
\label{thm:slln-clt}
For every $0<p<1$,
\begin{equation}
 \label{eq:slln}
 \frac{H_n}{n}\longrightarrow\mu_p=\frac{3a}{2+a}
 \qquad\text{almost surely}.
\end{equation}
Moreover,
\begin{equation}
 \label{eq:clt}
 \frac{H_n-n\mu_p}{\sqrt n}
 \xrightarrow{\mathrm d}\mathcal N(0,\sigma_p^2),
\end{equation}
where
\begin{equation}
 \label{eq:variance}
 \boxed{\quad
 \sigma_p^2=
 \frac{4a(1-a)(5-2a)}{(2+a)^3}>0.
 \quad}
\end{equation}
\end{theorem}

\begin{proof}
Let $g(i,\ell)\in\{-1,+1\}$ be the reward from state $i$ when the letter is
$\ell\in\{A,B\}$.  The conditional reward means are
\[
 f_N=q-p,\qquad f_Z=1,\qquad f_P=p-q.
\]
The finite chain is irreducible and aperiodic.  The ergodic theorem applied
to its transition rewards proves \eqref{eq:slln} with the stationary mean in
\cref{lem:stationary}.

To calculate fluctuations rather than cite an unspecified variance, solve
the Poisson equation
\[
 (I-P_p)h=f-\mu_p\one.
\]
Direct substitution gives the bounded solution
\begin{equation}
 \label{eq:poisson-solution}
 h_N=-\frac{2p}{1+p},\qquad h_Z=0,
 \qquad h_P=-\frac{2q}{1+q}.
\end{equation}
If $S_k$ is the projective state after $k$ letters, then
\begin{equation}
 \label{eq:martingale-difference}
 D_k=g(S_{k-1},X_k)-\mu_p+h(S_k)-h(S_{k-1})
\end{equation}
has conditional mean zero.  Indeed, conditioning on $S_{k-1}=i$ gives
$f_i-\mu_p+(P_ph)_i-h_i=0$.  The six state--letter values of $D_k$ are
\begin{center}
\renewcommand{\arraystretch}{1.13}
\begin{tabular}{c@{\qquad}cc}
\toprule
state & letter $A$ & letter $B$\\
\midrule
$N$ & $-2q/(1+q)$ & $2p/(1+q)$\\
$Z$ & $-2q(2p-1)/[(1+p)(1+q)]$
    & $2p(2p-1)/[(1+p)(1+q)]$\\
$P$ & $2q/(1+p)$ & $-2p/(1+p)$\\
\bottomrule
\end{tabular}
\end{center}
Consequently, the stationary conditional second moment is
\begin{align}
 \sum_i\pi_i\Exp(D_k^2\mid S_{k-1}=i)
 &=4a\left[
 \frac{\pi_N}{(1+q)^2}
 +\frac{\pi_Z(2p-1)^2}{(1+p)^2(1+q)^2}
 +\frac{\pi_P}{(1+p)^2}\right] \notag\\
 &=\frac{4a(1-a)(5-2a)}{(2+a)^3}.
 \label{eq:variance-derivation}
\end{align}
The last equality uses $(1+p)(1+q)=2+a$ and
$(2p-1)^2=1-4a$.

Summing \eqref{eq:martingale-difference} yields
\[
 H_n-n\mu_p=\sum_{k=1}^nD_k+h(S_0)-h(S_n).
\]
The final difference is bounded.  The conditional quadratic variation of
the martingale, divided by $n$, converges almost surely to
\eqref{eq:variance-derivation} by the finite-chain ergodic theorem.
The increments are bounded, so the conditional Lindeberg condition holds.
The martingale central limit theorem
\cite[Theorem~3.2]{HallHeyde1980} proves \eqref{eq:clt}.  Since
$0<a\leq1/4$, every factor in \eqref{eq:variance} is positive.
\end{proof}

The tilted cubic gives an additional algebraic check.  Once the literal
kernel is fixed, its variance calculation is independent of the Poisson
calculation above; the two proof routes as a whole are complementary.

\begin{proposition}[Perron derivative calculation]
\label{prop:perron-derivatives}
Let $\rho_p(t)=\pf(Q_p(e^t))$ for $0<p<1$.  Then
\[
 \rho_p(0)=1,\qquad
 (\log\rho_p)'(0)=\mu_p,\qquad
 (\log\rho_p)''(0)=\sigma_p^2.
\]
\end{proposition}

\begin{proof}
Set
\[
 F(r,t)=r^3+(2a-1-ae^{2t})r-ae^t.
\]
At $(r,t)=(1,0)$,
\[
 F_r=2+a,\quad F_t=-3a,\quad
 F_{rr}=6,\quad F_{rt}=-2a,\quad F_{tt}=-5a.
\]
Implicit differentiation of $F(\rho_p(t),t)=0$ gives
\[
 \rho_p'(0)=\frac{3a}{2+a}=\mu_p
\]
and
\[
 (2+a)\rho_p''(0)+6\mu_p^2-4a\mu_p-5a=0.
\]
Therefore
\begin{align*}
 (\log\rho_p)''(0)
 &=\rho_p''(0)-\mu_p^2\\
 &=\frac{5a+4a\mu_p-(8+a)\mu_p^2}{2+a}
 =\frac{4a(1-a)(5-2a)}{(2+a)^3}.
\end{align*}
This independently agrees with the martingale variance calculation in
\eqref{eq:variance-derivation}.
\end{proof}
